% !TeX program = pdflatex
% corpus/docs/calculo2_exercicios.tex — ticks C2-1–C2-21 (partes → Green).
%
% Prosa canónica; testemunhas em calculo2_resolve() (banco/conversa.c).
% Espinha: HIPÓTESES → DEFINIÇÃO → TRANSIÇÃO → LEI → TESTEMUNHA → CONCLUSÃO → VOLTA.
%
\documentclass[11pt,a4paper]{article}
\usepackage[utf8]{inputenc}
\usepackage[T1]{fontenc}
\usepackage[portuguese]{babel}
\usepackage{amsmath,amssymb,amsthm}
\usepackage[margin=2.4cm]{geometry}
\usepackage{booktabs}
\usepackage{xcolor}
\usepackage[hidelinks]{hyperref}

\newcommand{\code}[1]{\texttt{\small #1}}
\newcommand{\medido}[1]{\par\medskip\noindent{\small\color{gray}\textbf{Medido.} #1}\par\medskip}

\begin{document}
\title{Cálculo II --- exercícios\\
\large C2-1--C2-21: partes $\to$ Green}
\author{Tiffany --- corpus vivo}
\date{}
\maketitle

\begin{abstract}
\noindent Vinte e um exercícios (\code{C2\_21[]} em \code{banco/conversa.c}) ---
este ficheiro fecha o catálogo. Eixo: LOCAL $\to$ GLOBAL.
Operação ao vivo: \code{responde} / \code{calculo2 N}.
\end{abstract}

\section{como se le este documento}
Espinha dos sete slots. Lotes: C2-1--C2-8 (partes\ldots geométrica);
C2-9--C2-16 (raio\ldots cadeia multi); C2-17--C2-21 (Hessiana\ldots Green).
Colisões: \code{bolzano monotona}; \code{taylor momentos};
\code{lagrange multiplicadores}.

%======================================================================
\section{partes}\label{sec:partes-c2}
\textbf{Enunciado.} $\int u\,dv=uv-\int v\,du$, e sai DIRECTO da regra do produto.

\begin{enumerate}\itemsep2pt
\item \textbf{DEFINIÇÃO.} $u,v$ deriváveis num intervalo.
\item \textbf{TRADUÇÃO.} $\int u\,dv=uv-\int v\,du$.
\item \textbf{OPERAÇÃO.} Regra do produto lida ao contrário: $(uv)'=u'v+uv'$;
  integrar.
\item \textbf{LEI.} Teorema Fundamental --- partes e substituição são produto e
  cadeia invertidos.
\item \textbf{TESTEMUNHA.} Dois caminhos: $\int f g'$ vs $fg\big|_a^b-\int f'g$
  exactos.
\item \textbf{CONCLUSÃO.} A fórmula é a regra do produto integrada.
\item \textbf{VOLTA.} LOCAL $\to$ GLOBAL: produto local; integrar torna global
  (TFC).
\end{enumerate}
\medido{\code{calculo2\_resolve(1)}: $\int fg'=fg-\int f'g$ em $[0,2]$.}

%----------------------------------------------------------------------
\section{impropria}\label{sec:impropria}
\textbf{Enunciado.} $\int_1^\infty dx/x^p$ converge exactamente quando $p>1$.

\begin{enumerate}\itemsep2pt
\item \textbf{DEFINIÇÃO.} $p\ge1$ inteiro; integral (e série-p gémea).
\item \textbf{TRADUÇÃO.} Convergência exacta para $p>1$.
\item \textbf{OPERAÇÃO.} Não acumular soma parcial monstro --- comparar com soma
  que TELESCOPA.
\item \textbf{LEI.} Para $p\ge2$: $1/n^p\le1/(n(n-1))=1/(n-1)-1/n$; $\Sigma<1$.
\item \textbf{TESTEMUNHA.} Majoração vale em $p\ge2$, FALHA em $p=1$ (falhas
  contadas); harmónica por blocos.
\item \textbf{CONCLUSÃO.} Converge para $p>1$; diverge em $p=1$.
\item \textbf{VOLTA.} Objecto errado (soma estoura) $\neq$ guarda --- convergência
  decide-se por comparação.
\end{enumerate}
\medido{\code{calculo2\_resolve(2)} / \code{sr\_majora}: falha só em $p=1$.}

%----------------------------------------------------------------------
\section{sequencia}\label{sec:sequencia-c2}
\textbf{Enunciado.} $a_n=n/(n+1)\to1$, com o $N$ PROCURADO --- e $(-1)^n$ é o
gume.

\begin{enumerate}\itemsep2pt
\item \textbf{DEFINIÇÃO.} $a_n=n/(n+1)$.
\item \textbf{TRADUÇÃO.} $\forall\varepsilon>0\,\exists N:\ n>N\Rightarrow|a_n-1|<\varepsilon$.
\item \textbf{OPERAÇÃO.} $|n/(n+1)-1|=1/(n+1)$ exacto; achar $N$ em $\mathbb{Q}$.
\item \textbf{LEI.} $N$ não se promete --- PROCURA-SE (como o $\delta$ do Cálculo~I).
\item \textbf{TESTEMUNHA.} $N$ exibido por $\varepsilon$; gume $(-1)^n$ limitada e
  não monótona.
\item \textbf{CONCLUSÃO.} $n/(n+1)\to1$ com $N=\lceil1/\varepsilon\rceil$.
\item \textbf{VOLTA.} Escada de observadores (\code{thm:escada}): quem não
  estabiliza não converge.
\end{enumerate}
\medido{\code{calculo2\_resolve(3)} / \code{sq\_acha\_N}: $N$ para cada
$\varepsilon$.}

%----------------------------------------------------------------------
\section{bolzano monotona}\label{sec:bolzano-c2}
\textbf{Enunciado.} Limitada $+$ monótona $\Rightarrow$ convergente; e só limitada
NÃO basta.

\begin{enumerate}\itemsep2pt
\item \textbf{DEFINIÇÃO.} $(a_n)$ limitada.
\item \textbf{TRADUÇÃO.} Limitada $+$ monótona $\Rightarrow$ convergente.
\item \textbf{OPERAÇÃO.} Completude: crescente limitada tem $\sup$ e converge para
  ele; monotonia impede oscilação.
\item \textbf{LEI.} As duas hipóteses operam; retirada a monotonia, cai.
\item \textbf{TESTEMUNHA.} Gume $(-1)^n$: limitada, NÃO monótona, NÃO converge ---
  dois limites de subsucessão.
\item \textbf{CONCLUSÃO.} Limitada sozinha NÃO chega; monotonia é operacional.
\item \textbf{VOLTA.} Bolzano--Weierstrass salva subsucessão --- aqui podem ser
  duas, com limites distintos.
\end{enumerate}
\medido{\code{calculo2\_resolve(4)}: $N$ da sequencia + gume $(-1)^n$.}

%----------------------------------------------------------------------
\section{serie}\label{sec:serie-c2}
\textbf{Enunciado.} A série é o objecto; o VALOR é que precisa do limite.

\begin{enumerate}\itemsep2pt
\item \textbf{DEFINIÇÃO.} $\sum a_n x^n$ com coeficientes racionais.
\item \textbf{TRADUÇÃO.} $S_N=\sum_{n=0}^{N}a_n$; converge $\Leftrightarrow S_N\to S$.
\item \textbf{OPERAÇÃO.} Objecto $=$ vector de coeficientes (soma, produto, deriva
  sem limite); valor precisa do limite.
\item \textbf{LEI.} Série FORMAL --- como Dirichlet pelos coeficientes
  (\code{dirichlet.h}).
\item \textbf{TESTEMUNHA.} Identidade formal $(1-x)\sum x^n=1$ com resíduo zero;
  somas parciais à parte.
\item \textbf{CONCLUSÃO.} Objecto existe antes; convergência é pergunta sobre o
  valor.
\item \textbf{VOLTA.} LOCAL $\to$ GLOBAL: coeficientes locais; convergência global;
  extremos à parte.
\end{enumerate}
\medido{\code{calculo2\_resolve(5)}: objecto formal vs valor.}

%----------------------------------------------------------------------
\section{termo geral}\label{sec:termo-geral}
\textbf{Enunciado.} $a_n\to0$ é NECESSÁRIO e não suficiente --- a harmónica é o
gume.

\begin{enumerate}\itemsep2pt
\item \textbf{DEFINIÇÃO.} $\sum a_n$; tipicamente $a_n=1/n^p$.
\item \textbf{TRADUÇÃO.} $\sum a_n$ converge $\Rightarrow a_n\to0$ --- recíproca
  FALSA.
\item \textbf{OPERAÇÃO.} Termo geral LOCAL; convergência GLOBAL --- o salto não é
  automático.
\item \textbf{LEI.} Necessário $\neq$ suficiente.
\item \textbf{TESTEMUNHA.} Em $p=1$ o termo vai a zero e a série DIVERGE (blocos
  $\ge1/2$).
\item \textbf{CONCLUSÃO.} $a_n\to0$ necessário e não suficiente; harmónica prova.
\item \textbf{VOLTA.} LOCAL $\to$ GLOBAL na forma mais crua deste andar.
\end{enumerate}
\medido{\code{calculo2\_resolve(6)} / \code{sr\_termo\_a\_zero}: gume harmónica.}

%----------------------------------------------------------------------
\section{serie p}\label{sec:serie-p}
\textbf{Enunciado.} $\sum1/n^p$: a comparação que TELESCOPA, e a soma exacta
ESTOURA.

\begin{enumerate}\itemsep2pt
\item \textbf{DEFINIÇÃO.} $\sum_{n\ge1}1/n^p$.
\item \textbf{TRADUÇÃO.} Converge exactamente quando $p>1$.
\item \textbf{OPERAÇÃO.} Objecto da soma parcial exacta ESTOURA o \texttt{long};
  comparação telescópica fica pequena.
\item \textbf{LEI.} $1/n^p\le1/(n-1)-1/n$ para $p\ge2$; $\Sigma\le1-1/N<1$.
\item \textbf{TESTEMUNHA.} Telescopagem $=$ forma fechada; majoração $0$ falhas em
  $p=2,3$; falha em $p=1$.
\item \textbf{CONCLUSÃO.} Converge $p\ge2$; diverge $p=1$ --- sem monstro.
\item \textbf{VOLTA.} Objecto errado (soma) $\neq$ guarda; critério estrutural.
\end{enumerate}
\medido{\code{calculo2\_resolve(7)} / \code{sr\_telescopa}: identidade exacta.}

%----------------------------------------------------------------------
\section{geometrica}\label{sec:geometrica}
\textbf{Enunciado.} $(1-x)\cdot\sum x^n=1$ como série FORMAL, sem convergência
nenhuma.

\begin{enumerate}\itemsep2pt
\item \textbf{DEFINIÇÃO.} Série de potências $\sum x^n$.
\item \textbf{TRADUÇÃO.} $(1-x)\sum_{n\ge0}x^n=1$.
\item \textbf{OPERAÇÃO.} Produto de Cauchy exacto: coeficientes cancelam-se aos
  pares; sobra $1$ --- sem convergência.
\item \textbf{LEI.} Identidade entre vectores de coeficientes.
\item \textbf{TESTEMUNHA.} $a_0=1$ e zero nulos acima; somas parciais em $x=1/2$ à
  parte.
\item \textbf{CONCLUSÃO.} Identidade FORMAL; convergência é pergunta SEPARADA.
\item \textbf{VOLTA.} Extremos do raio ($x=\pm1$) examinam-se à parte --- o salto
  local/global não é sozinho.
\end{enumerate}
\medido{\code{calculo2\_resolve(8)} / \code{sr\_geometrica}: resíduo zero.}

%----------------------------------------------------------------------
\section{raio}\label{sec:raio}
\textbf{Enunciado.} $|x-c|<R$ converge, $>R$ diverge, e os extremos à parte.

\begin{enumerate}\itemsep2pt
\item \textbf{DEFINIÇÃO.} $\sum a_n(x-c)^n$ com coeficientes racionais.
\item \textbf{TRADUÇÃO.} $|x-c|<R$ converge; $|x-c|>R$ diverge.
\item \textbf{OPERAÇÃO.} Raio decide o aberto; extremos ($|x-c|=R$) examinados
  SEPARADAMENTE.
\item \textbf{LEI.} Série formal vs valor --- o raio é pergunta sobre o valor.
\item \textbf{TESTEMUNHA.} $\sum x^n/n$ com $R=1$: em $x=1$ diverge (harmónica);
  em $x=-1$ converge (alternada).
\item \textbf{CONCLUSÃO.} Interior do disco pelo raio; bordo à parte.
\item \textbf{VOLTA.} LOCAL $\to$ GLOBAL: extremos são onde o salto não se faz
  sozinho.
\end{enumerate}
\medido{\code{calculo2\_resolve(9)}: extremos $x=\pm1$ separados.}

%----------------------------------------------------------------------
\section{taylor momentos}\label{sec:taylor-c2}
\textbf{Enunciado.} $f$ pelos seus coeficientes --- e são os MOMENTOS do
observador.

\begin{enumerate}\itemsep2pt
\item \textbf{DEFINIÇÃO.} $f$ suficientemente derivável em torno de $a$.
\item \textbf{TRADUÇÃO.} $f(x)=\sum_{n=0}^{N}f^{(n)}(a)/n!\,(x-a)^n+R_N(x)$.
\item \textbf{OPERAÇÃO.} Ler a função pelos COEFICIENTES --- objectos algébricos
  FINITOS; grau $=$ quanto se vê.
\item \textbf{LEI.} Identidades de derivação $=$ identidades de coeficientes.
\item \textbf{TESTEMUNHA.} $\exp'=\exp$, $\sin'=\cos$, $\cos'=-\sin$ termo a termo;
  $\sin^2+\cos^2=1$ como série.
\item \textbf{CONCLUSÃO.} Taylor lê por coeficientes; identidades sobrevivem.
\item \textbf{VOLTA.} Coeficientes $=$ momentos $\Phi_m$ (\code{thm:escada});
  truncamento $=$ degrau do observador (\code{thm:teto-oito}).
\end{enumerate}
\medido{\code{calculo2\_resolve(10)}: identidades nos coeficientes; $0$ falhas.}

%----------------------------------------------------------------------
\section{exp seno cosseno}\label{sec:exp-seno-cosseno}
\textbf{Enunciado.} $\exp'=\exp$, $\sin'=\cos$, $\cos'=-\sin$: nos COEFICIENTES.

\begin{enumerate}\itemsep2pt
\item \textbf{DEFINIÇÃO.} Séries $e^x$, $\sin x$, $\cos x$ truncadas.
\item \textbf{TRADUÇÃO.} Derivação termo a termo das três.
\item \textbf{OPERAÇÃO.} Verificar coeficiente a coeficiente, sem avaliar.
\item \textbf{LEI.} Mesma lei de Taylor --- identidades de coeficientes.
\item \textbf{TESTEMUNHA.} Três derivações $0$ falhas; $\sin^2+\cos^2=1$ com
  $a_0=1$ e nulos acima.
\item \textbf{CONCLUSÃO.} As três identidades fecham nos coeficientes.
\item \textbf{VOLTA.} Pitágoras sem ângulo e sem decimal --- só série formal.
\end{enumerate}
\medido{\code{calculo2\_resolve(11)} / \code{sr\_exp},\code{sr\_sin},\code{sr\_cos}.}

%----------------------------------------------------------------------
\section{parciais}\label{sec:parciais}
\textbf{Enunciado.} $\partial f/\partial x$ e $\partial f/\partial y$ formais e
exactas.

\begin{enumerate}\itemsep2pt
\item \textbf{DEFINIÇÃO.} $f(x,y)$ polinomial com coeficientes racionais.
\item \textbf{TRADUÇÃO.} Parciais formais e exactas.
\item \textbf{OPERAÇÃO.} Derivada formal num índice; o outro congelado --- Cálculo~I
  numa família de índices.
\item \textbf{LEI.} Mesma operação do quociente / regra formal.
\item \textbf{TESTEMUNHA.} Gradiente exibido; Schwarz em 625 polinómios.
\item \textbf{CONCLUSÃO.} Parciais exactas em polinómios.
\item \textbf{VOLTA.} Par dual: $df$ funcional vs $\nabla f$ vector
  (\code{thm:produto-dual}).
\end{enumerate}
\medido{\code{calculo2\_resolve(12)}: parciais formais; varredura Schwarz.}

%----------------------------------------------------------------------
\section{gradiente}\label{sec:gradiente-c2}
\textbf{Enunciado.} $\nabla f$ é o DUAL: o funcional $df$ lido como vector.

\begin{enumerate}\itemsep2pt
\item \textbf{DEFINIÇÃO.} Mesmo $f$ polinomial.
\item \textbf{TRADUÇÃO.} $\nabla f=(\partial f/\partial x,\partial f/\partial y)$.
\item \textbf{OPERAÇÃO.} $\nabla f$ VECTOR; $df$ FUNCIONAL --- métrica identifica.
\item \textbf{LEI.} Sem produto interno há $df$ e não há $\nabla f$.
\item \textbf{TESTEMUNHA.} $\nabla f$ em ponto com coordenadas exactas.
\item \textbf{CONCLUSÃO.} $\nabla f$ $=$ par ordenado das parciais.
\item \textbf{VOLTA.} Nomear um sem o outro é meio teorema
  (\code{thm:produto-dual}).
\end{enumerate}
\medido{\code{calculo2\_resolve(13)} / \code{p2\_gradiente}: vector exacto.}

%----------------------------------------------------------------------
\section{schwarz}\label{sec:schwarz}
\textbf{Enunciado.} $f_{xy}=f_{yx}$ --- VERIFICADO nos coeficientes, não assumido.

\begin{enumerate}\itemsep2pt
\item \textbf{DEFINIÇÃO.} $f(x,y)$ polinomial.
\item \textbf{TRADUÇÃO.} $f_{xy}=f_{yx}$.
\item \textbf{OPERAÇÃO.} Derivar em $x$ depois $y$ multiplica por $i$ e por $j$;
  ordem não importa em $\mathbb{Q}$.
\item \textbf{LEI.} Comutatividade da multiplicação --- Schwarz sem continuidade das
  segundas.
\item \textbf{TESTEMUNHA.} $625$ polinómios: $0$ falhas na ordem das parciais.
\item \textbf{CONCLUSÃO.} Parciais comutam para polinómios.
\item \textbf{VOLTA.} Identidade verificável --- não postulado.
\end{enumerate}
\medido{\code{calculo2\_resolve(14)}: $625$ polinómios; $0$ falhas.}

%----------------------------------------------------------------------
\section{derivada linear}\label{sec:derivada-linear}
\textbf{Enunciado.} $Df(a)$ NÃO é um número: é um OPERADOR LINEAR local.

\begin{enumerate}\itemsep2pt
\item \textbf{DEFINIÇÃO.} $f:\mathbb{R}^n\to\mathbb{R}^m$ derivável.
\item \textbf{TRADUÇÃO.} $Df(a):\mathbb{R}^n\to\mathbb{R}^m$ --- não um número.
\item \textbf{OPERAÇÃO.} Salto do andar: Cálculo~I dava inclinação (número); aqui
  MATRIZ $=$ melhor aproximação linear local.
\item \textbf{LEI.} Único $L$ linear com $f(a+h)=f(a)+Lh+o(h)$; em polinómios
  $o(h)$ exacto (grau $\ge2$).
\item \textbf{TESTEMUNHA.} Jacobiana de aplicação linear; $\det=$ factor de volume.
\item \textbf{CONCLUSÃO.} Objecto muda de tipo --- liga cálculo à álgebra linear.
\item \textbf{VOLTA.} Matriz local; volume que move é global
  ($\Lambda^n T=\det T$).
\end{enumerate}
\medido{\code{calculo2\_resolve(15)} / \code{p2\_jacobiana}: matriz e $\det$.}

%----------------------------------------------------------------------
\section{cadeia multi}\label{sec:cadeia-multi}
\textbf{Enunciado.} $D(f\circ g)=Df(g(x))\cdot Dg(x)$: compor vira MULTIPLICAR
matrizes.

\begin{enumerate}\itemsep2pt
\item \textbf{DEFINIÇÃO.} $f,g$ deriváveis entre espaços.
\item \textbf{TRADUÇÃO.} $D(f\circ g)(x)=Df(g(x))\,Dg(x)$.
\item \textbf{OPERAÇÃO.} Compor funções $\Leftrightarrow$ multiplicar operadores; a
  ordem importa.
\item \textbf{LEI.} Linearidade da aproximação; $\det(AB)=\det A\cdot\det B$.
\item \textbf{TESTEMUNHA.} Dois caminhos: jacobiana da composta vs produto das
  jacobianas --- quatro entradas.
\item \textbf{CONCLUSÃO.} Compor É multiplicar; volumes compõem-se.
\item \textbf{VOLTA.} Teorema do gato: VOLUME como segunda vida --- dualidade /
  medida.
\end{enumerate}
\medido{\code{calculo2\_resolve(16)}: $Df\cdot Dg=$ jacobiana da composta.}

%----------------------------------------------------------------------
\section{hessiana}\label{sec:hessiana}
\textbf{Enunciado.} $H$ simétrica, e a assinatura decide: mínimo, máximo ou SELA.

\begin{enumerate}\itemsep2pt
\item \textbf{DEFINIÇÃO.} $a$ ponto crítico: $\nabla f(a)=0$.
\item \textbf{TRADUÇÃO.} $H>0\Rightarrow$ mínimo; $H<0\Rightarrow$ máximo;
  indefinida $\Rightarrow$ SELA.
\item \textbf{OPERAÇÃO.} Segunda ordem: $f(a+h)-f(a)=\frac12 h^T H h+\ldots$; o
  SINAL da forma decide --- a ASSINATURA, não a matriz olhada.
\item \textbf{LEI.} Formas quadráticas: procurar TESTEMUNHAS, não ler a matriz.
\item \textbf{TESTEMUNHA.} $x^2-y^2$ sela (par de vectores $+/-$); $x^2+y^2$ mín;
  $-x^2-y^2$ máx; $x^2$ DEGENERADA (critério não decide).
\item \textbf{CONCLUSÃO.} Assinatura classifica; degenerado é quarto estado.
\item \textbf{VOLTA.} Hessiana LOCAL classifica um ponto; extremo GLOBAL precisa do
  que o Cálculo~I já ensinou.
\end{enumerate}
\medido{\code{calculo2\_resolve(17)} / \code{p2\_classifica}: quatro estados.}

%----------------------------------------------------------------------
\section{lagrange multiplicadores}\label{sec:lagrange-c2}
\textbf{Enunciado.} $\nabla f=\lambda\nabla g$ no extremo --- e o gume é a
regularidade do vínculo.

\begin{enumerate}\itemsep2pt
\item \textbf{DEFINIÇÃO.} Extremizar $f(x,y)$ sujeito a $g(x,y)=c$.
\item \textbf{TRADUÇÃO.} $\nabla f=\lambda\nabla g$.
\item \textbf{OPERAÇÃO.} No extremo vinculado, $\nabla f$ perpendicular à curva de
  nível de $g$ --- logo paralelo a $\nabla g$.
\item \textbf{LEI.} Mesma caça de testemunhas das formas.
\item \textbf{TESTEMUNHA.} Gume: $\nabla g=0$ na origem em $g=x^3-y^2$; nenhum
  $\lambda$ --- método não se aplica a singularidade do vínculo.
\item \textbf{CONCLUSÃO.} Condição com regularidade $\nabla g\neq0$; sem ela, cai.
\item \textbf{VOLTA.} Gume pedido pelo eval: regularidade do vínculo é hipótese
  operacional.
\end{enumerate}
\medido{\code{calculo2\_resolve(18)}: gume $\nabla g=(0,0)$; sem $\lambda$.}

%----------------------------------------------------------------------
\section{fubini}\label{sec:fubini-c2}
\textbf{Enunciado.} As duas ordens dão o mesmo, por caminhos intermédios DISTINTOS.

\begin{enumerate}\itemsep2pt
\item \textbf{DEFINIÇÃO.} $f$ polinomial no rectângulo $R$.
\item \textbf{TRADUÇÃO.} $\iint_R f\,dA=\int(\int f\,dy)\,dx=\int(\int f\,dx)\,dy$.
\item \textbf{OPERAÇÃO.} Dois objectos intermédios DIFERENTES (polinómio em $x$ vs
  em $y$); só o número final coincide --- teorema, não tautologia.
\item \textbf{LEI.} INTERIOR $\leftrightarrow$ BORDA na forma fraca (duas iterações).
\item \textbf{TESTEMUNHA.} Duas ordens exactas iguais em $[0,2]\times[0,3]$.
\item \textbf{CONCLUSÃO.} Fubini vale por caminhos distintos.
\item \textbf{VOLTA.} LOCAL $\to$ GLOBAL: o número é global; os caminhos são locais
  distintos.
\end{enumerate}
\medido{\code{calculo2\_resolve(19)}: $\int\!dy\,dx=\int\!dx\,dy$.}

%----------------------------------------------------------------------
\section{jacobiano}\label{sec:jacobiano-c2}
\textbf{Enunciado.} $dV=|\det J|\,du\,dv$ --- o determinante como MEDIDA DE VOLUME.

\begin{enumerate}\itemsep2pt
\item \textbf{DEFINIÇÃO.} Mudança $(u,v)\mapsto(x,y)$.
\item \textbf{TRADUÇÃO.} $dV=|\det J|\,du\,dv$.
\item \textbf{OPERAÇÃO.} $|\det J|$ não é correcção: é a MEDIDA --- acção do
  operador no volume (andar exterior).
\item \textbf{LEI.} $\det$ como $\Lambda^n T$.
\item \textbf{TESTEMUNHA.} Jacobiana de $(u,v)\mapsto(2u+v,u+3v)$: $\det$ exacto.
\item \textbf{CONCLUSÃO.} Determinante $=$ factor de volume.
\item \textbf{VOLTA.} Matriz local; volume global --- mesma frase da cadeia multi.
\end{enumerate}
\medido{\code{calculo2\_resolve(20)}: $\det J$ factor de volume.}

%----------------------------------------------------------------------
\section{green}\label{sec:green-c2}
\textbf{Enunciado.} $\oint_{\partial D}=\iint_D$: o INTERIOR contra a BORDA,
medido dos dois lados.

\begin{enumerate}\itemsep2pt
\item \textbf{DEFINIÇÃO.} $R=[x_0,x_1]\times[y_0,y_1]$; campo $(P,Q)$ polinomial.
\item \textbf{TRADUÇÃO.} $\oint_{\partial R}P\,dx+Q\,dy=\iint_R(Q_x-P_y)\,dA$.
\item \textbf{OPERAÇÃO.} Esquerda $=$ soma na BORDA (quatro segmentos); direita $=$
  soma no INTERIOR --- contas diferentes, mesmo número.
\item \textbf{LEI.} INTERIOR $\leftrightarrow$ BORDA --- forma forte de LOCAL
  $\to$ GLOBAL.
\item \textbf{TESTEMUNHA.} $P=-y$, $Q=x$: borda $=$ área $=12$; $625$ campos,
  $0$ divergências.
\item \textbf{CONCLUSÃO.} Green liga interior à borda --- global pelo contorno.
\item \textbf{VOLTA.} \code{thm:borda-dirac}; adjunto $\delta\dashv\varepsilon$;
  Green/Gauss/Stokes $=$ mesma frase em dimensões diferentes.
\end{enumerate}
\medido{\code{calculo2\_resolve(21)} / \code{p2\_green\_*}: $625$ campos; $0$ falhas.}

%======================================================================
\section{o que fica no conversa}
A prosa canónica C2-1--C2-21 é este ficheiro. \code{calculo2\_resolve} \textbf{mede}.
Catálogo de Cálculo II fechado.

\end{document}
